\documentclass[border=8pt]{standalone}
\usepackage{tikz}
\usepackage{xcolor}
\usetikzlibrary{arrows.meta, positioning, fit, backgrounds, calc}
\usepackage[T1]{fontenc}
\usepackage[utf8]{inputenc}
\usepackage{kotex}

\definecolor{Ink}{RGB}{26,34,43}
\definecolor{Deep}{RGB}{0,57,127}
\definecolor{Faint}{RGB}{124,138,153}
\definecolor{Match}{RGB}{215,110,55}

\begin{document}
\begin{tikzpicture}[
    tt/.style     = {font=\ttfamily\small, color=Ink, inner sep=1pt},
    mm/.style     = {font=\ttfamily\small, color=Ink, inner xsep=2pt, inner ysep=1.5pt},
    match/.style  = {draw=Match, thick, rounded corners=2pt},
    rulel/.style  = {draw=Ink, line width=0.4pt},
    arr/.style    = {-{Latex[length=6pt]}, Faint, thick},
    lbl/.style    = {font=\sffamily\small, color=Deep},
    capT/.style   = {font=\sffamily\footnotesize, color=Faint},
    x=1cm, y=1cm,
    node distance=0pt,
]

%% =================== TREE A (LEFT) ===================
\node[lbl] (labA) at (-6.0, 5.15) {증명나무 A};
\node[tt]  (A_top) at (-6.0, 4.5) {\{\} $\vdash$ 1 $\Rightarrow$ 1};
\draw[rulel] ($(A_top.south west)+(-0.20,-0.08)$) -- ($(A_top.south east)+(0.20,-0.08)$);
% Bottom line split so we can highlight the output memory only
\node[tt] (A_bot_pre) at (-7.15, 3.85) {\{\} $\vdash$ x := 1;\ $\Rightarrow$};
\node[mm, match, right=1pt of A_bot_pre] (A_mem) {\{x: 1\}};

%% =================== TREE B (RIGHT) ===================
\node[lbl] (labB) at (5.3, 6.15) {증명나무 B};
\node[tt] (B_topL) at ( 3.2, 5.5) {\{x: 1\} $\vdash$ x $\Rightarrow$ 1};
\node[tt] (B_topR) at ( 7.4, 5.5) {\{x: 1\} $\vdash$ 2 $\Rightarrow$ 2};
\draw[rulel] ($(B_topL.south west)+(-0.20,-0.08)$) -- ($(B_topR.south east)+(0.20,-0.08)$);
\node[tt] (B_mid) at (5.3, 4.85) {\{x: 1\} $\vdash$ x + 2 $\Rightarrow$ 3};
\draw[rulel] ($(B_mid.south west)+(-0.20,-0.08)$) -- ($(B_mid.south east)+(0.20,-0.08)$);
% Bottom line: highlight the leading input memory
\node[mm, match] (B_mem) at (2.0, 4.2) {\{x: 1\}};
\node[tt, right=1pt of B_mem] (B_rest) {$\vdash$ y := x + 2;\ $\Rightarrow$ \{x: 1, y: 3\}};

%% =================== MERGED TREE (BOTTOM) ===================
% Left premise (split to highlight output memory)
\node[tt] (M1_pre) at (-4.6, 0.95) {\{\} $\vdash$ x := 1;\ $\Rightarrow$};
\node[mm, match, right=1pt of M1_pre] (M1_mem) {\{x: 1\}};

% Right premise (split to highlight input memory)
\node[mm, match] (M2_mem) at (1.7, 0.95) {\{x: 1\}};
\node[tt, right=1pt of M2_mem] (M2_rest) {$\vdash$ y := x + 2;\ $\Rightarrow$ \{x: 1, y: 3\}};

\draw[rulel] ($(M1_pre.south west)+(-0.20,-0.15)$) -- ($(M2_rest.south east)+(0.20,-0.15)$);
\node[tt] (M_bot) at ($(M1_pre.south east)!0.5!(M2_rest.south west) + (0,-0.55)$)
   {\{\} $\vdash$ x := 1;\ y := x + 2;\ $\Rightarrow$ \{x: 1, y: 3\}};

\node[lbl] at ($(M_bot.south)+(0,-0.4)$) {합쳐진 증명나무};

%% =================== CONNECTING ARROWS ===================
\draw[arr] ($(A_mem.south)+(0,-0.05)$) to[out=-90, in=90]
           ($(M1_mem.north)+(0,0.05)$);
\draw[arr] ($(B_mem.south)+(0,-0.05)$) to[out=-90, in=90]
           ($(M2_mem.north)+(0,0.05)$);

%% =================== SIDE CAPTION ===================
\node[capT, align=left, anchor=west] at (9.6, 3.0)
  {{\color{Match}\bfseries 연결점 메모리 일치}\\
   A의 결과 메모리 $=$ B의 시작 메모리\\
   \ttfamily\{x: 1\}\normalfont\sffamily\footnotesize\color{Faint}로 같아야\\
   두 나무를 Seq 규칙으로\\
   이을 수 있다.};

\end{tikzpicture}
\end{document}
